\clearpage
\section{Reference-equality successor coherence and exact costs}
\label{sec:copying-equality-coherence}
Both reference-equality branches now have complete successor-state coherence
for the actual combined machine. Two independently chosen contextual matches
at the same comparison request produce recursively isomorphic whole graphs,
and the isomorphism preserves the resumed token. This closes both comparison
cases of Appendix~B.1's machine obligation. It does not close the whole
machine package. Section~\ref{sec:copying-assignment-coherence} subsequently
completes assignment, and Section~\ref{sec:copying-dereference-coherence}
completes dereference. Section~\ref{sec:copying-arithmetic-coherence} completes
arithmetic. Sections~\ref{sec:copying-application-coherence} and \ref{sec:copying-successor-coherence} complete application and copying successor comparison; refocusing and the specimen laws remain open.

The operational rule is still the paper's Figure~19(d), discussed in
\href{https://arxiv.org/html/1907.01257v5\#S10.SS1}{HTML Section~10.1}.
The two outcomes depend on whether the argument instances lead to the same
store edge. The new proof compares the complete reducts, extending the earlier
uniqueness of that Boolean decision. All original edges outside the comparison
head are retained, including instances, contraction trees, stores and their
attached contents. The replacement introduced below is a proof construction;
it is not a new rule or a change to the language.

The dependency chain is: unique source head and store decision; explicit
local replacement isomorphisms; state-level substitution congruences from
Section~\ref{sec:copying-allocation-coherence}; induction over actual rewrites;
bounded and unbounded successor coherence; exact successful-cost laws.
Neither Proposition~12.2's terms nor Figure~16's translation changes.
The four broad remaining packages stay open. The cautious estimate remains
50--65\% of mathematical effort, with low confidence: closing these two
cases does not establish the dependent template argument.

\Needspace{12\baselineskip}
\subsection{A result with two explicit weakenings}

The finite-index helpers say that a one-element type contains only zero,
and that every element of the three-element type is zero, one or two.
Their proofs use the bounds carried by each finite index. They support
exhaustive incidence checks and inverse-map proofs below.

\lean{Copying/DiscardResult}{9}{9}{DiscardResult.fin_one}

\lean{Copying/DiscardResult}{10}{10}{DiscardResult.fin_three}

\termdef{discard-result-graph}
For a constant $c$, the \emph{result-and-discard graph} $D_c$ has interface
$\star\to\star\otimes\star$, no internal vertices and three edges. Edge zero
is the constant $c$, with the sole input as its source and no targets.
Edges one and two are term weakenings, with no sources and with the first
and second output, respectively, as their targets. The label, source and
target functions specify precisely these three cases. The graph therefore
returns a value while explicitly retaining the two old argument wires.

\lean{Copying/DiscardResult}{13}{14}{DiscardResult.label}

\lean{Copying/DiscardResult}{16}{16}{DiscardResult.source}

\lean{Copying/DiscardResult}{17}{18}{DiscardResult.target}

Enumerating the three edges proves well-formedness. Incidence lists are
empty or singleton; only the constant uses the input and each weakening
uses its own output. This gives uniqueness and coverage. Labels have exactly
the specified interfaces, proving typing. Both carriers are finite.
Primitive relabelling interprets $D_c$ at any depth allowance and preserves
these properties.

\lean{Copying/DiscardResult}{21}{73}{DiscardResult.base}

\lean{Copying/DiscardResult}{76}{77}{DiscardResult.net}

\Needspace{12\baselineskip}
\subsection{The canonical equality successor}

\termdef{canonical-equality-successor}
Suppose $e$ in $G$ is labelled by atom equality, $v$ is the request wire, and
$b$ is a Boolean. The \emph{canonical equality successor} is
\[
 E(G,e,v,b)=R(G,e,D_{\mathrm{bool}(b)},v),
\]
where $R$ is the replacement state of
Section~\ref{sec:copying-allocation-coherence}. The filler first transports
$D_{\mathrm{bool}(b)}$ along the interface equalities implied by the head label.
Then $E$ replaces only $e$, resumes search at the retained image of $v$, and
keeps every other edge occurrence. Defining this graph requires no termination
or stability assumption about the stored contents. Its agreement with actual
rewriting is a separate theorem below.

\lean{Copying/EqualityReplacement}{9}{14}{AtomEquality.filler}

\lean{Copying/EqualityReplacement}{17}{19}{AtomEquality.canonical}

A recursive isomorphism $f:G\cong H$ carries the canonical successor to
$E(H,f(e),f(v),b)$. Primitive labels are preserved by recursive isomorphism,
so the renamed head has the same operation. Apply the proved replacement
renaming construction. The two boundary casts on the filler compose to the
single cast required in $H$; proof irrelevance identifies their equality
witnesses. The state isomorphism transports both the graph and token.

\lean{Copying/EqualityReplacement}{22}{34}{AtomEquality.canonical_rename}

\Needspace{12\baselineskip}
\subsection{The same-store local rule}

The local head is the unique edge in the comparison component. Its label
is atom equality and its source is the request input, proving the head
predicate. The abbreviation \texttt{replaced} names the graph component of
$E(G,e,v,\mathrm{true})$ for this local source, with an arbitrary matched
contraction tree. No canonical presentation of that tree is required.

\lean{Copying/EqualityReplacement}{38}{38}{AtomEquality.Same.head}

\lean{Copying/EqualityReplacement}{39}{41}{AtomEquality.Same.head_at}

\lean{Copying/EqualityReplacement}{44}{45}{AtomEquality.Same.replaced}

The forward internal-vertex map retains each frame vertex and each
contraction-tree vertex, while removing the extra sum tags introduced by
one-edge plugging. The inverse restores these tags. Primitive corollas and
$D_c$ have empty internal carriers, so all other cases are impossible.
Case analysis proves the two composites are identities, yielding the
internal-vertex bijection.

\lean{Copying/EqualityReplacement}{48}{55}{AtomEquality.Same.vertexTo}

\lean{Copying/EqualityReplacement}{57}{65}{AtomEquality.Same.vertexFrom}

\lean{Copying/EqualityReplacement}{67}{78}{AtomEquality.Same.vertices}

For edges, retain each left-instance, right-instance, tree and store
occurrence. The old comparison edge is excluded by the subtype used for
retained host edges. Send the three new edges to the result, left weakening
and right weakening components of the original target. The inverse sends
those components back to indices zero, one and two and restores every
retained edge. Exhausting the finite indices proves both inverse laws;
in the excluded-head case the retained-edge inequality gives a contradiction.

\lean{Copying/EqualityReplacement}{81}{88}{AtomEquality.Same.edgeTo}

\lean{Copying/EqualityReplacement}{90}{97}{AtomEquality.Same.edgeFrom}

\lean{Copying/EqualityReplacement}{99}{111}{AtomEquality.Same.edges}

The frame-port equation follows by splitting a port into input, internal
vertex or output. The tree-port equation then splits the tree port: an
internal port follows directly from the vertex map, while input and output
ports reduce to the frame equation at the corresponding tree attachment.
This checks the whole interface of an arbitrary tree, not only its root.

\lean{Copying/EqualityReplacement}{113}{117}{AtomEquality.Same.map_frame}

\lean{Copying/EqualityReplacement}{120}{127}{AtomEquality.Same.map_tree}

Assemble the two bijections into a graph isomorphism. Internal wire labels
and retained edge labels agree by construction. Primitive edge incidences
reduce directly; tree incidences are lists of inserted ports, so mapping each
entry uses the tree-port equation. The three fresh edges have exactly the
sources and targets of the target result and weakenings. Finally, extend the
graph isomorphism to a state isomorphism: both tokens have search mode and
the same external input. This proves canonicality of the actual same-store
local rule, with its complete retained tree and store.

\lean{Copying/EqualityReplacement}{130}{163}{AtomEquality.Same.replacement_iso}

\lean{Copying/EqualityReplacement}{166}{168}{AtomEquality.Same.canonical_local}

\Needspace{12\baselineskip}
\subsection{The distinct-store local rule}

The distinct-store construction follows the same method with two arbitrary
trees and two distinct store components. Its unique comparison head again
satisfies the head predicate. Here \texttt{replaced} is the graph component
of $E(G,e,v,\mathrm{false})$. These definitions select the original local
rule's head and Boolean; they do not merge stores with equal contents.

\lean{Copying/EqualityDifferentReplacement}{8}{8}{AtomEquality.Different.head}

\lean{Copying/EqualityDifferentReplacement}{9}{11}{AtomEquality.Different.head_at}

\lean{Copying/EqualityDifferentReplacement}{14}{15}{AtomEquality.Different.replaced}

The vertex maps retain frame vertices and the separately tagged vertices
of both trees. All other internal cases are empty. Restoring and deleting
these tags proves the inverse laws. The edge maps retain both instances,
both trees and both store occurrences. Only the comparison head is replaced;
the three new indices map to result and argument weakenings. Exhaustive
finite-index cases prove the edge maps inverse, while the retained-head
case is excluded by its inequality proof.

\lean{Copying/EqualityDifferentReplacement}{18}{27}{AtomEquality.Different.vertexTo}

\lean{Copying/EqualityDifferentReplacement}{29}{39}{AtomEquality.Different.vertexFrom}

\lean{Copying/EqualityDifferentReplacement}{41}{52}{AtomEquality.Different.vertices}

\lean{Copying/EqualityDifferentReplacement}{55}{64}{AtomEquality.Different.edgeTo}

\lean{Copying/EqualityDifferentReplacement}{66}{75}{AtomEquality.Different.edgeFrom}

\lean{Copying/EqualityDifferentReplacement}{77}{89}{AtomEquality.Different.edges}

The frame-port equation again fixes inputs, outputs and retained frame
vertices. Separate left-tree and right-tree equations reduce boundary ports
to that frame equation and internal ports to their own tagged vertex map.
Thus neither tree can be confused with the other, even when their labels
or presentations coincide.

\lean{Copying/EqualityDifferentReplacement}{91}{95}{AtomEquality.Different.map_frame}

\lean{Copying/EqualityDifferentReplacement}{98}{105}{AtomEquality.Different.map_tree}

\lean{Copying/EqualityDifferentReplacement}{108}{115}{AtomEquality.Different.map_right_tree}

The full graph isomorphism checks labels and incidence lists using both
tree-port equations, and checks the three replacement edges directly.
It extends to the actual successor state because the resumed input token
is fixed. The statement compares complete successors, so retaining the
second store and all its attachments is part of the conclusion.

\lean{Copying/EqualityDifferentReplacement}{118}{161}{AtomEquality.Different.replacement_iso}

\lean{Copying/EqualityDifferentReplacement}{164}{166}{AtomEquality.Different.canonical_local}

\Needspace{12\baselineskip}
\subsection{Independence of every actual contextual match}

An actual rewrite at an atom-equality head supplies a store-identity
decision. Apply the already proved recognition theorem to that rewrite;
source-head uniqueness forces its recognised operation to be atom equality.
The other operation constructors, including each arithmetic code, are
incompatible. The remaining comparison recognition contains the required
Boolean decision. No decision premise is added to the final coherence theorem.

\lean{Copying/EqualityCoherence}{9}{17}{Combined.rewrite_equality_decision}

To reflect a comparison into the changed component of a context, first
reflect the unique head using the actual inner rewrite and its match
embedding. Obtain an inner decision from that rewrite, transport it to the
host, and compare it with the given host decision. Uniqueness of store
resolution makes the Booleans equal. Thus the reflected head, its edge image
and the decision are all justified without assuming globally injective
boundary attachment.

\lean{Copying/EqualityCoherence}{20}{28}{Combined.rewrite_equality_reflect}

For a local rewrite, compare its distinguished head with the specified
head. All non-comparison constructors contradict the operation label.
In the same-store and distinct-store constructors, head uniqueness identifies
the edge and decision uniqueness identifies the Boolean. The explicit local
state isomorphism from the preceding subsections then applies.

\lean{Copying/EqualityCoherence}{31}{67}{Combined.local_equality_canonical}

Induct over the actual contextual rewrite derivation. The local case is
settled above. In the renaming case, reflect the head and decision, apply the
induction hypothesis, and compose with canonical renaming and the supplied
renaming of the actual output. For plugging, reflect into the filler, use the
induction hypothesis there, and compose replacement associativity with the
lifted state isomorphism. For simultaneous expansion, reflect into the selected
component; replacement/expansion interchange and equality of the other
components yield the complete result. Every constructor of the actual rewrite
relation is covered.

\lean{Copying/EqualityCoherence}{70}{93}{Combined.rewrite_equality_canonical}

\termdef{equality-successor-coherence}
\emph{Equality successor coherence} says that any two actual rewrites from
the same graph and request vertex, at an atom-equality head, have recursively
isomorphic successor states. Derive the Boolean from the first rewrite;
canonicality sends both successors to the same $E(G,e,v,b)$. Invert one
isomorphism and compose with the other. This allows independently chosen
hosts, tree presentations, context decompositions and output representations.
It assumes neither machine determinism nor agreement of those choices.

\lean{Copying/EqualityCoherence}{96}{107}{Combined.equality_rewrite_deterministic}

In the full step relation, a source in rewrite mode cannot take a search
step. Eliminating those cases reduces step coherence to rewrite coherence.
For the unbounded allowance-indexed machine, each step lifts a bounded step
at the same allowance, so bounded coherence gives represented successor
states. The existing transport of executions along representation then gives
successor equivalence for every cost preorder, including exact equality.
This is equivalence of the successor executions; it is not yet the contextual
equivalence asserted by Proposition~12.2.

\lean{Copying/EqualityCoherence}{110}{119}{Combined.equality_step_deterministic}

\lean{Copying/EqualityCoherence}{129}{138}{Combined.equality_unbounded_deterministic}

\lean{Copying/EqualityCoherence}{141}{144}{Combined.equality_unbounded_equiv}

\Needspace{12\baselineskip}
\subsection{Exact costs and arbitrary hosts}

For an admitted comparison step $s\to t$, successor coherence gives
\[
 \mathsf{TerminatesIn}(n,s)
 \quad\Longleftrightarrow\quad
 \exists m,\ n=m+1\ \land\ \mathsf{TerminatesIn}(m,t).
\]
A zero-step successful run would make $s$ final, contradicting its step.
For a positive run, compare its first successor with $t$ by coherence and
transport the remaining run at the same cost. Conversely, prepend the
chosen comparison step. The generic cost lemma of the allocation milestone
encapsulates this argument.

\lean{Copying/EqualityCoherence}{122}{126}{Combined.equality_cost_spectrum}

Serial attachment of a graph to a continuation preserves every actual
search interaction, with the token carried by the left insertion map.
This follows from the proved simultaneous-expansion search law, selecting
the left component of the serial pattern.

\lean{Copying/EqualityCost}{9}{19}{Search.serial_left}

Attach arbitrary store contents to either comparison fragment. Map its
head into the serial composition and lift the original local compute step.
Full step coherence shows that every competing successor is isomorphic to
this explicit successor. More generally, place the fragment in an arbitrary
typed host edge, with equality casts for the matching interfaces. Head
transport and contextual step transport give the same conclusion. There is
no hypothesis that the inserted input remains an external input, so the host
may resume execution after the comparison returns.

\lean{Copying/EqualityCost}{22}{28}{AtomEquality.Same.contents_successor}

\lean{Copying/EqualityCost}{75}{82}{AtomEquality.Different.contents_successor}

\lean{Copying/EqualityCost}{60}{72}{AtomEquality.Same.host_successor}

\lean{Copying/EqualityCost}{113}{123}{AtomEquality.Different.host_successor}

For serially attached contents, the only successful comparison cost is
seven. There are five preparatory search interactions: enter the left
argument, return from its instance, move to the right argument, return from
its instance, and request comparison. There is then one compute step and
one return from the Boolean constant. The final input is external in this
serial composition. Apply search-cost inversion five times, the comparison
cost law once, and search-cost inversion once more; the final state has
only zero successful cost. Finite arithmetic gives exactly seven and also
constructs the run of that length. No stability, termination or acyclicity
assumption about the stored contents is used. The statement excludes every
alternative successful cost under the full combined relation, rather than
merely exhibiting a seven-step local run.

\lean{Copying/EqualityCost}{31}{57}{AtomEquality.Same.contents_exact_cost}

\lean{Copying/EqualityCost}{85}{110}{AtomEquality.Different.contents_exact_cost}

\Needspace{12\baselineskip}
\subsection{Unchanged translated contents, nested boxes and padding}

Successful costs in the unbounded machine are precisely the successful
costs of a fixed bounded representative. Reflect an unbounded finite run to
its original allowance, preserving the endpoint and finality; conversely,
lift a bounded run. This general lemma is used in both regression families.

\lean{Copying/EqualityCoherenceExamples}{9}{21}{Combined.successful_unbounded_iff}

The first closed regression graph attaches the unchanged translation of
an arbitrary closed term $t$ to a same-store comparison with a noncanonical
binary atom tree. The second attaches two unchanged copies of that same
translation to two distinct stores. It uses a tensor of the translated
contents and one canonical one-input atom tree per store. Serial and tensor
label-preservation laws, together with the existing translation theorems,
prove focus freedom and hole freedom, so these are valid closed programs.
The comparison redexes are explicitly constructed local graph scenarios;
no claim identifies them with the whole translation of a surface comparison
expression.

\lean{Copying/EqualityCoherenceExamples}{24}{29}{EqualityCoherenceExamples.sharedContents}

\lean{Copying/EqualityCoherenceExamples}{32}{41}{EqualityCoherenceExamples.distinctContents}

Each regression has exactly the successful cost seven in the unbounded
machine after any finite amount of depth padding. First use exact padding
reflection, then the bounded/unbounded successful-cost correspondence, then
the appropriate arbitrary-contents theorem. The serial input map fixes the
external input, giving the original initialization state. Separately, map
the local store decision into the serial graph: the shared-store result is
true, and the distinct-store result is false despite identical contents.

\lean{Copying/EqualityCoherenceExamples}{44}{51}{EqualityCoherenceExamples.shared_exact_cost}

\lean{Copying/EqualityCoherenceExamples}{54}{61}{EqualityCoherenceExamples.distinct_exact_cost}

\lean{Copying/EqualityCoherenceExamples}{64}{69}{EqualityCoherenceExamples.shared_decision}

\lean{Copying/EqualityCoherenceExamples}{72}{80}{EqualityCoherenceExamples.distinct_decision}

Specialise both cost results to the unchanged translation of
$\Omega=(\lambda x.xx)(\lambda x.xx)$. Comparison terminates without
executing the stored application or opening its boxes; this does not prove
$\Omega$'s divergence. The arbitrary-term results also cover deeper nested boxes.

\lean{Copying/EqualityCoherenceExamples}{83}{88}{EqualityCoherenceExamples.omega_contents}

